\chapter{Grundlagen der Ereignisbehandlung}
\renewcommand{\chaptertitle}{Grundlagen der Ereignisbehandlung}

\lehead[]{\normalfont\sffamily\hspace*{-2.00cm}\textcolor{white}{\colorbox{lightblue}{\makebox[1.60cm][r]{\thechapter}}}\hspace{0.17cm}\textcolor{lightblue}{\chaptertitle}}
\rohead[]{\textcolor{lightblue}{\chaptertitle}\normalfont\sffamily\hspace*{0.17cm}\textcolor{white}{\colorbox{lightblue}{\makebox[1.60cm][l]{\thechapter}}}\hspace{-2.00cm}}
%\chead[]{}
\rehead[]{\textcolor{lightblue}{AvHG, Inf, My}}
\lohead[]{\textcolor{lightblue}{AvHG, Inf, My}}

Bisher liefen unsere Programme nach einem festen Plan ab: Das Fenster wurde
geöffnet, \lstinline|myPaint()| zeichnete Figuren, und ein \lstinline|Timer|
sorgte im Takt für Bewegung. Moderne Programme mit Buttons, Textfeldern und
Menüs funktionieren jedoch anders -- sie sind \textbf{ereignisgesteuert}
(\emph{event-driven}). Das Programm wartet ab, bis der Benutzer etwas tut, und
reagiert dann gezielt darauf.


\section{Die vier Beteiligten an einem Ereignis}

Wenn du in einem Fenster auf einen Button klickst oder eine Taste drückst,
arbeiten vier Instanzen im Hintergrund zusammen:

\begin{compactenum}[1.]
  \item \textbf{Der Benutzer (bzw.\ die Hardware):}\\
  Durch einen Mausklick oder Tastendruck sendet die Hardware ein elektrisches
  Signal an den Computer.

  \item \textbf{Das Betriebssystem (OS):}\\
  Das Betriebssystem nimmt das Signal entgegen und ermittelt, welches Programm
  gerade den \emph{Fokus} besitzt (welches Fenster also aktiv im Vordergrund
  liegt bzw.\ unter dem Mauszeiger liegt). Nur an dieses Programm leitet das
  Betriebssystem die Nachricht weiter -- alle anderen offenen Programme bleiben
  unbehelligt.

  \item \textbf{Die Laufzeitumgebung (JVM / Swing):}\\
  Die Java Virtual Machine empfängt die Nachricht vom Betriebssystem, verpackt
  alle Details (z.\,B.\ Mauskoordinaten, gedrückte Taste oder angeklickter
  Button) in ein \textbf{Event-Objekt} (z.\,B.\ ein \myClass{ActionEvent}) und
  stellt es in eine Warteschlange (\emph{Event Queue}). Ein spezieller Thread
  -- der \emph{Event Dispatch Thread (EDT)} -- arbeitet diese Warteschlange
  der Reihe nach ab.

  \item \textbf{Unser Programm (der Listener):}\\
  Die Komponente (z.\,B.\ ein \myClass{JButton}) benachrichtigt alle
  Objekte, die sich zuvor als \glqq Zuhörer\grqq\ (\emph{Listener}) bei ihr
  angemeldet haben, und ruft dort die passende Reaktionsmethode auf.
\end{compactenum}

\medskip
\begin{center}
\fbox{\parbox{0.92\textwidth}{\centering\small
\textbf{Hardware} (Klick) $\longrightarrow$
\textbf{Betriebssystem} (prüft aktives Fenster) $\longrightarrow$
\textbf{JVM / EDT} (erzeugt Event-Objekt) $\longrightarrow$
\textbf{Listener} (führt deinen Code aus)
}}
\end{center}


\section{Das Listener-Prinzip in Java}

Woher weiß ein Button, \emph{welcher} Code ausgeführt werden soll, wenn er
angeklickt wird? Die Programmierer von Java konnten schließlich nicht ahnen, was
dein Button einmal tun soll.

Die Lösung heißt \textbf{Delegation (Listener-Prinzip)}:
\begin{compactitem}
  \item Für verschiedene Arten von Ereignissen gibt es passende
  \glqq Zuhörer\grqq -Typen (\emph{Listener-Interfaces}) -- für einen
  Button-Klick ist das der \myClass{ActionListener} (für Tastatur- oder
  Mausbewegungen gäbe es entsprechend \myClass{KeyListener} oder
  \myClass{MouseListener}).
  \item Jeder Button verwaltet intern eine Liste seiner angemeldeten
  \myClass{ActionListener} (theoretisch könnten also mehrere Programmteile
  gleichzeitig auf denselben Klick reagieren -- in unseren Programmen melden
  wir meist genau einen Zuhörer an).
  \item Über die Methode \lstinline|addActionListener(...)| meldet man beim
  Button an, welcher Code beim Klick ausgeführt werden soll. Sobald der Button
  angeklickt wird, ruft er automatisch die Methode \linebreak
  \lstinline|actionPerformed(ActionEvent e)| bei allen angemeldeten
  Zuhörern auf.
\end{compactitem}

\subsection{Kurzschreibweise mit Lambda-Ausdrücken}

In modernem Java (und im \emph{WindowBuilder}, den wir im nächsten Kapitel
benutzen) schreibt man diese Anmeldung sehr kompakt mithilfe eines sogenannten
\emph{Lambda-Ausdrucks} (\lstinline|e -> ...|):

\begin{lstlisting}
btnBerechnen.addActionListener(e -> {
    // Dieser Code wird immer dann ausgeführt,
    // wenn der Benutzer auf den Button klickt!
    berechneErgebnis();
});
\end{lstlisting}

Du musst diese Zeile meistens nicht einmal selbst tippen: Wenn du im
WindowBuilder einen Doppelklick auf einen Button machst, erzeugt Eclipse
genau dieses Gerüst automatisch für dich. Deine Aufgabe besteht dann nur noch
darin, die gewünschte Reaktion zwischen die geschweiften Klammern zu
schreiben.
